\documentclass[10pt,a4paper]{amsart}
\usepackage[margin=19mm]{geometry}
\usepackage{amsmath,amssymb,mathtools}
\usepackage[scr=rsfs,cal=euler]{mathalfa}
\usepackage{xfrac}
\usepackage[unicode,hidelinks,bookmarks=false]{hyperref}
\newcommand{\ie}{i.e.\ }
\newcommand{\cf}{cf.\ }
\newcommand{\resp}{respectively\ }
\newcommand{\Subsection}{\subsection}
\providecommand{\nobreakdash}{\nobreak\hbox{-}}
\setlength{\emergencystretch}{2em}
\multlinegap0pt
\usepackage{bm}
\usepackage{bbm}
\usepackage{dsfont}
\usepackage{xspace}
\usepackage{cancel}
\usepackage{mathtools}
\usepackage{amsthm}
\makeatletter
\@ifundefined{theorem}{\newtheorem{theorem}{Theorem}}{}
\makeatother
\newcommand{\tr}{\operatorname{tr}}
\title[Resummed effective actions and heat kernels: the Worldline a]{Resummed effective actions and heat kernels: the Worldline approach and Yukawa assisted pair creation}
\author{Filippo Fecit, Sebasti\'an Franchino-Vi\~nas, Francisco D. Mazzitelli}
\date{Source: arXiv:2501.17094v2 (CC BY 4.0)}
\begin{document}
\maketitle

\noindent\emph{Source and licence.} Resummed effective actions and heat kernels: the Worldline approach and Yukawa assisted pair creation. Version arXiv:2501.17094v2 (CC BY 4.0), distributed under the Creative Commons Attribution 4.0 International licence, as recorded in the arXiv licence metadata for this exact version and shown on the abstract page https://arxiv.org/abs/2501.17094v2. Full rights notice: licenses/arxiv-2501.17094-CC-BY-4.0.txt.\par\smallskip

\noindent\textit{Editorial scope.} Counted unit: the Theorem whose source label is \texttt{th:GY\_periodic} in the article (source0.tex lines 1247--1254, source label \texttt{th:GY\_periodic}), delivered together with the article's own complete original proof of it (source0.tex lines 1257--1292, one \texttt{proof} environment of 36 lines). No mathematics is written, repaired or re-derived: statement and proof are the article's own text line for line. Required context retained uncounted: eq:Delta\_SI (lines 444--446); eq:weighted\_derivatives (lines 1111--1113); BCY (lines 1121--1123); GY (lines 1148--1150). Every definition, display and label of those lines is retained unchanged. Omitted from this package and not redistributed: 1--443, 447--1110, 1114--1120, 1124--1147, 1151--1246, 1255--1256, 1293--1338. Nothing inside a retained excerpt is omitted or altered: every retained body is delivered exactly as the article's lines are written, so no omission receipt of any kind accompanies this package. Citations: the retained excerpts carry no citation command at all, so no bibliography entry of the article is transcribed and no reference is redistributed. Reference adaptation: none was needed and none was made; every reference inside the delivered unit resolves inside it, so no unresolved reference or citation is produced. Printed slips kept literal: no printed word, symbol, hypothesis or number of the counted statement or of its proof was altered. Layout and class: the article's own class is \texttt{article} with packages including jheppub, amsthm, cancel, empheq, float, blindtext, enumitem, amsmath, latexsym, amssymb, amsfonts, bm, bbm, mathtools; the wrapper is the lane's pinned \texttt{amsart} editorial wrapper at 10pt on A4, which restates the theorem-like environments the retained text needs and adds the article's own macro definitions that the retained text needs (\texttt{\textbackslash tr}), with extra wrapper packages (bm, bbm, dsfont, xspace, cancel, mathtools). The delivered numbering is the unit's own. Assets: no retained span contains a figure environment, a graphics-inclusion command, a PDF-page inclusion, a table or a TikZ picture; no image file is copied into this package, the package declares no asset dependency and no image file is redistributed with it. Licence: version arXiv:2501.17094v2 of this article is distributed under the Creative Commons Attribution 4.0 International licence, as recorded in the arXiv licence metadata for this exact version and shown on the abstract page https://arxiv.org/abs/2501.17094v2. A full rights notice is delivered at licenses/arxiv-2501.17094-CC-BY-4.0.txt. Characters: every retained excerpt is pure ASCII, so the delivered engine, XeLaTeX with its default Latin Modern text font, reports no missing character.
\begin{equation} \label{eq:Delta_SI}
{\Delta^{\mu\nu}_{\rm SI} (t,t')=-\frac12 \delta^{\mu\nu}\partial^2_t\delta(t-t')+2 \, \left(\Omega^2\right)^{\mu\nu}(\bar{x})\delta(t-t')}\ .
\end{equation}

\begin{align}\label{eq:weighted_derivatives}
 \mathbf{v}^{(j)}(z):=P_i(z)\frac{{\rm d}}{{\rm d}z}\mathbf{u}^{(j)}(z)\ ;
\end{align}

\begin{equation} \label{BCY}
Y(0)=\mathbb{1}_{2r \times 2r}\ . 
\end{equation}

\begin{equation} \label{GY}
\frac{\mathrm{Det}(L_1)}{\mathrm{Det}(L_2)}=\frac{\mathrm{det}(M+N \, Y_1(T))}{\mathrm{det}(M+N \, Y_2(T))}\ .
\end{equation}

\begin{theorem}\label{th:GY_periodic}
Consider an operator of the form given in Eq.~\eqref{eq:Delta_SI}, defined on the domain of periodic functions. Be $u_a^{(b)}$ the $a$th component of the $b$th solution of the homogeneous equation $L_1 \mathbf{u}^{(b)}=0$ and $\mathbf{v}^{(b)}$ its corresponding weighted derivative, as in Eq.~\eqref{eq:weighted_derivatives}. 
Then, the relevant determinant for the GY theorem satisfies
\begin{align}\label{eq:theorem}
\mathrm{det}\Big(M+N \, Y_{\Delta, \mathrm{PBC}}(T)\Big)
=\det \big(2-2\cosh(2\Omega T)\big)\ .
\end{align}
\end{theorem}

\begin{proof}
Note first that, in contrast to the DBC, to compute the functional determinant using the generalized GY theorem we do need the whole set of $2r$ solutions of the homogeneous equation
\begin{equation}
\Delta_{\rm SI}^{\mu\nu} \, \Phi_\nu^{(\Pi)}=0\ , \quad \Phi^{(\Pi)}_\mu=\begin{cases}
 \phi_\mu^{(\rho)} \quad \text{if} \quad \Pi=1,\dots,r \\
 \psi_\mu^{(\rho)} \quad \text{if} \quad \Pi=r+1,\dots,2r
\end{cases}\ .
\end{equation}
 Considering Eq.~\eqref{BCY} as initial conditions, namely 
\begin{align}
\phi_\mu^{(\rho)}(0)&=\delta^\rho_\mu\ , \quad \frac{{\rm d}\phi}{{\rm d}z}_\mu^{(\rho)}(0)=\mathbb{0}\ ,
\\
\psi_\mu^{(\rho)}(0)&=\mathbb{0}\ , \quad \frac{{\rm d}\psi}{{\rm d}z}_\mu^{(\rho)}(0)=\delta^\rho_\mu\ ,
\end{align}
an explicit computation gives
\begin{equation}
\phi_\mu^{(\rho)}(z)=\dot{\psi}_\mu^{(\rho)}(z)=\Big[\cosh(2\Omega z)\Big]^\rho{}\Big._\mu\ .
\end{equation}
Note then that in arbitrary dimensions we can straightforwardly prove that the $n$th power of the matrix $Y_\Delta$ is
\begin{align}
 Y^n_{\Delta, \mathrm{PBC}} (T)= \begin{pmatrix}
 \cosh(2n\Omega T) & \sinh(2n\Omega T)
 \\
 \sinh(2n\Omega T) & \cosh(2n\Omega T)
 \end{pmatrix},
\end{align}
and consequently, from Eq.~\eqref{GY}, we have the equalities
\begin{align}
 \begin{split}
 \det(\mathbb{1}_{2r} - Y_{\Delta, \mathrm{PBC}}) 
 &= \mathrm{e}^{\tr \log [2-2 \cosh(2\Omega T)]}
 = \det \big(2-2\cosh(2\Omega T)\big) \ ,
 \end{split}
\end{align}
which indeed agree with the RHS of Eq.~\eqref{eq:theorem}. 
\end{proof}

\end{document}
